\documentclass[12pt,a4paper]{article}
\usepackage[utf8]{inputenc}
\usepackage{amsmath, amssymb, amsfonts}
\usepackage{geometry}
\geometry{margin=2.5cm}
\usepackage{hyperref}
\usepackage{pgfplots}
\usepackage{fancyhdr}
\usepackage{listing}
\usepackage{listings}
\usepackage{titlesec}
\titleformat{\section}
  {\normalfont\Large\bfseries\centering}  % Format: font size, bold, centered
  {\thesection}{1em}{}
\pagestyle{fancy}
\pgfplotsset{compat=1.18}
\fancyhead{}
\fancyfoot{}
\title{Laboratorio 2}
\author{}
\date{}
\fancyhead[L]{Gabriele Perna}          % Left side
\fancyhead[C]{Laboratorio 2}     % Center
\fancyhead[R]{2026/2027}             % Right side
\fancyfoot[C]{\thepage}           % Page number in the center of footer
\let\oldsection\section
\renewcommand{\section}{\clearpage\oldsection}

\begin{document}

\maketitle
\tableofcontents

\section{Introduzione}
In questo corso lavoreremo con programmi vicini al linguaggio macchina.\\
Lavoreremo inoltre con:
\begin{itemize}
    \item UNIX
    \item thread
    \item processi
    \item programmazione di sistema 
\end{itemize}
\paragraph{Codice Teams:} r407rxv

\section{Il terminale UNIX}
Il \textbf{terminale} è una interfaccia testuale per dialogare con il sistema.\\
Dentro il terminale gira una \textbf{shell}, ovvero il programma che interpreta i comandi.\\
Il terminale mostra il \textbf{prompt}, riceve i comandi e visualizza l'output.
\subsection{Comandi}
\subsubsection{Struttura base comandi Linux}
Esiste una struttura base per molti comandi, ovvero questa:
\begin{center}
    comando + opzione + argomento + (opzione + argomento)
\end{center}
per esempio:
\begin{lstlisting}
    $ gcc -Wall hello.c -o hello
\end{lstlisting}
\subsubsection{--help e man}
\begin{lstlisting}
    $ gcc --help
    $ man gcc
\end{lstlisting}
Entrambi questi comandi aiutano a comprendere specifici comandi in modi diversi, --help è un aiuto generico mentre "man" è il manuale.
\subsubsection{pwd}
\begin{lstlisting}
    $ pwd
\end{lstlisting}
Il simbolo $\$$ indica che la shell è pronta a ricevere un comando (potrebbe non essere $\$$ a volte).\\
\textbf{pwd} è il comando "print working directory" e mostra la posizione della directory corrente all'interno del \textbf{file system}.\\
Il file system inizia dalla \textbf{root} ("$/$") e si articola nella directory, formando un \textbf{path}, ovvero un percorso attraverso directory, dove ogni passaggio è separato da "$/$" (potrebbe variare in altri ambienti).\\
Esistono due tipi di path:
\begin{itemize}
    \item \textbf{Assoluto:} parte dalla root.
    \item \textbf{Relativo:} parte dalla directory corrente.
\end{itemize}
\subsubsection{ls}
\begin{lstlisting}
    $ ls
\end{lstlisting}
\textbf{ls} elenca tutto il contenuto della directory corrente, directory e file.\\
Può essere usata con alcuni \textbf{attributi}, aggiungibili con un "dash" (-) come in questo caso:
\begin{lstlisting}
    $ ls -a
\end{lstlisting}
Ecco gli attributi:
\begin{itemize}
    \item \textbf{a}: include tutti gli elementi nascosti
    \item \textbf{l}: include informazioni avanzate (data, grandezza...)
\end{itemize}
\subsubsection{cd}
\begin{lstlisting}
    $ cd esercizi
    $ cd /users/gabri/esercizi
\end{lstlisting}
\textbf{cd} significa "change directory" e si sposta nella directory specificata da un path relativo o assoluto.\\
Percorsi speciali:
\begin{itemize}
    \item \textbf{Directory corrente}: specificato da un punto, ha lo stesso valore di pwd.
    \begin{lstlisting}
       $ cd . 
    \end{lstlisting}
    \item \textbf{Directory padre}: specificato da un doppio punto.
    \begin{lstlisting}
       $ cd .. 
    \end{lstlisting}
\end{itemize}
\subsubsection{mkdir}
\begin{lstlisting}
    $ mkdir esercizi
\end{lstlisting}
\textbf{mkdir} significa "make directory" e crea una directory del nome specificato nella cartella corrente.
\section{Linguaggio C}
Il linguaggio \textbf{C} è un linguaggio compilato, pertanto, prima di eseguire un file bisogna convertirlo in binario.\\
Esistono diversi standard C: C89, C95, C99, C11, C17, C23.\\
\subsection{Compilatore}
Il compilatore più conosciuto e usato per C è \textbf{gcc}, ovvero "GNU C Compiler", anche se adesso viene definito come "GNU Compiler Collection" perche diventato un progetto più ampio che include anche altri compilatori come quello di Fortran e C++. 
\begin{lstlisting}
    gcc hello.c -o hello
\end{lstlisting}
In questo caso, il comando:
\begin{itemize}
    \item compila il file "hello.c"
    \item rinomina l'eseguibile come "hello"
\end{itemize}
Ogni programma eseguito ritorna un risultato per determinare se è terminato con successo o meno, si ottiene cosi:
\begin{lstlisting}
    $ gcc hello.c
    $ hello.c
    $ echo $?
    0
\end{lstlisting}
\subsubsection{Warnings}
I \textbf{warnings} sono importanti e non vanno ignorati, per compilare e stamparli tutti si fa cosi:
\begin{lstlisting}
    $ gcc -Wall hello.c
\end{lstlisting}
Ogni warning contiene un messaggio e la location (riga e carattere) di esso.
\subsection{Librerie}
Nei programmi in C è possibile utilizzare delle \textbf{direttive di preprocessazione} (o del preprocessore) che permettono di importare funzioni e strumenti utili allo svolgimento del programma:
\begin{lstlisting}
    #include <stdio.h>
\end{lstlisting}
\subsection{Variabili}
In C possiamo usare le \textbf{variabili}. Si possono inizializzare ad un valore di un tipo specifico, assegnato al momento della \textbf{dichiarazione}.
\begin{lstlisting}
    int x;
    x = 32;
\end{lstlisting}
\subsection{printf}
Questo è un comando base per stampare qualcosa su console:
\begin{lstlisting}
    printf("Hello World");
\end{lstlisting}
\subsubsection{Placeholders}
Per stampare le variabili bisogna utilizzare i \textbf{placeholders}:
\begin{lstlisting}
    int x = 10;
    printf("%d",x)
\end{lstlisting}
Bisogna scegliere il placeholder corretto per riuscire a stampare una variabile; la scelta va in base al tipo:
\begin{itemize}
    \item \textbf{\%d}: numeri interi (int)
    \item \textbf{\%i}: numeri interi (int)
    \item \textbf{\%f}: numeri decimali (double)
    \begin{itemize}
        \item Per stampare uno specifico numero di cifre decimali si aggiunge un "." dopo "\%" e poi si specifica il numero:
        \begin{lstlisting}
            printf(%.2f,x)
        \end{lstlisting}
        Questa stampa mostrerà per esempio solo 2 cifre decimali. Usato sugli altri placeholder questo provoca l'ingrandimento di padding a destra (numero negativo) o sinistra (numero positivo) di ciò che viene stampato.
    \end{itemize}
    \item \textbf{\%c}: caratteri (char)
    \item \textbf{\%s}: stringa
    \item \textbf{\%u}: numeri interi senza segno (unsigned int)
    \item \textbf{\%x}: numero intero in esadecimale
    \item \textbf{\%p}: indirizzo di memoria, puntatore
\end{itemize}
\subsection{Commenti}
I commenti in C possono essere inseriti in vari modi:
\begin{itemize}
    \item Commenti su una riga: aperti da "//", non vengono chiusi
    \item Commenti su più righe: delimitati da "/*" e da "*/".
\end{itemize}
\end{document}